% !TeX program = xelatex
\documentclass[UTF8,scheme=plain,a4paper,11pt]{ctexart}
\usepackage[margin=25mm,top=21mm,bottom=22mm]{geometry}
\usepackage{amsmath,amssymb,amsthm,booktabs}
\usepackage[unicode,hidelinks]{hyperref}
\hypersetup{pdftitle={三元域双参数前缀主子式的最坏代价}}
\ctexset{section={format=\large\bfseries\raggedright,name={},number=\arabic{section},aftername=\quad}}
\newcommand{\F}{\mathbb F_3}
\newtheorem{thm}{定理}
\newtheorem{lem}[thm]{引理}
\renewcommand{\proofname}{证明}
\renewcommand{\refname}{参考文献}
\setlength{\emergencystretch}{1em}
\begin{document}
\fontsize{12}{18.5}\selectfont
\begin{center}
{\LARGE\bfseries 三元域双参数前缀主子式的最坏代价\par}
\medskip 研究短稿 \quad 2026 年 11 月 21 日
\end{center}
\noindent\textbf{摘要}\quad 对归一化的三元域双参数矩阵族，考察使全部参数在同一行序下通过真前缀主子式检验所需的最少行加法次数。我们证明，$m\ge3$ 时其最大值为 $\lceil2m/3\rceil$，且两列斜率线性无关的矩阵族能达到该值。上界来自一个短块选取引理，所选块移除后仍有非零标量总和。下界来自任意成功操作序列的连通分量数。

\section{模型与六个小规模算例}
取 $u,b\in\F^m$ 和 $m\times2$ 矩阵 $W$，令 $z(a)=b+Wa$，参数 $a$ 遍历 $\F^2$。归一化条件为 $u^\top W=0$ 和 $u^\top b=1$。记每行的斜率电荷为 $p_i=u_iW_{i,*}\in\F^2$，标量电荷为 $s_i=u_i b_i\in\F$。对行指标集 $B$，有
\[
 \det\bigl(I-u_B z_B(a)^\top\bigr)=1-\sum_{i\in B}(s_i+p_i a).
\]
若存在一个行序，使每个真前缀（前 $k$ 行，$1\le k<m$）的这一行列式对所有 $a$ 均非零，就称该状态通过检验。行变换 $P$ 同时作用为 $(Pu,P^{-\top}W,P^{-\top}b)$。置换和逐行变号免费，每次 $R_i\gets R_i+cR_j$ 收费一步，其中 $i\ne j$，$c=\pm1$。记通过检验所需的最少收费步数为 $C$，其在全部归一化输入上的最大值为 $M_2(m)$。这里允许两列斜率线性相关。
\begin{thm}\label{th}
对每个正整数 $m$，
\[
 M_2(m)=\min\{m-1,\lceil2m/3\rceil\}.
\]
前两项为 $M_2(1)=0$、$M_2(2)=1$。当 $m\ge3$ 时，$M_2(m)=\lceil2m/3\rceil$，且存在 $\operatorname{rank}W=2$ 的输入达到最大值。
\end{thm}
把相差免费操作的输入归为一类，从所有通过检验的类同时出发，沿可逆行加法作广度优先搜索，可得下表。三列分秩给出最坏代价，横线表示该秩没有输入。
\begin{center}\small
\begin{tabular}{r r r r r r}\toprule
行数 $m$ & 免费操作等价类数 & 秩零 & 秩一 & 秩二 & $M_2(m)$\\\midrule
1&1&0&--&--&0\\2&28&0&1&--&1\\3&436&0&2&2&2\\
4&4,935&1&2&3&3\\5&44,898&1&3&4&4\\6&345,726&1&4&4&4\\\bottomrule
\end{tabular}
\end{center}
按各类代表的带标号输入数加权，可核对每一秩的输入总数。归一化保证 $u\ne0$，两列斜率又同属 $u^\perp$，故秩二至少需要三行。

\newpage
\section{移除短块后保留非零标量总和}
下文的块是非空行指标集。$\F^2$ 中，三个不同的点共线当且仅当它们的和为零。任意五个不同的点必含一条完整的线。否则每个方向的平行线至多贡献两对点，四个方向至多贡献八对，但五点共有十对。
\begin{lem}\label{short}
给定至少十二对 $(p_i,s_i)\in\F^2\times\F$ 和 $t\in\F^\times$，可以选出至多三行组成块 $B$，使 $\sum_Bp_i=0$ 且 $\sum_Bs_i\ne t$。
\end{lem}
\begin{proof}
把标量乘以 $t^{-1}$，可设 $t=1$。假设每个至多三行且斜率和为零的块，其标量和都为一。记出现的斜率集合为 $P$。同一斜率至多出现三次。若它出现四次，则任取其中三行的标量和都为一。比较这些三行组，四个标量必须相等，而三个相等标量之和为零。

若 $x\ne0$ 且 $-x\in P$，把斜率为 $x$ 的各行分别与同一个 $-x$ 行配对，可知它们的标量相同。若 $x$ 在 $P$ 的某条完整线上，固定该线另外两点的代表行，也有同样的结论。这两种情况下，斜率 $x$ 都至多出现两次，其标量记为 $f(x)$。

零斜率行的标量只能为一，故零斜率至多出现一次。若它出现，$P$ 就不能同时含 $x,-x$（$x\ne0$），因为在这对行后加上零斜率行，标量和仍须为一。若没有这样的相反对，则至多有四个非零斜率。恰有四个时，把原点添入这四点，所得五点中的完整线必由三个非零点组成。这三种斜率各至多出现两次。即使计入一行零斜率，总行数也至多 $1+3\cdot2+3=10$。非零斜率至多三种时，上界同为十行。

若恰有一对相反斜率，此外至多还有三种斜率。若另有至多两种，总行数至多十。若另有三种，这五点含一条完整线，它至少经过这三种斜率中的两种。因此至少四种斜率各至多出现两次，总行数至多十一。

三对相反斜率不可能同时出现。选取符号后可将它们记为 $\pm x,\pm y,\pm(x+y)$。由三条配对等式，六个标量总和为零。将 $x,y,-x-y$ 所在直线和它的相反直线对应的两条等式相加，却得到总和为二。

最后设有两对相反斜率 $\pm x,\pm y$。若另外两个方向均有斜率出现，通过变号或交换 $x,y$，可设 $w=x+y$ 和 $z=x-y$ 都出现。写 $A=f(x)$、$B=f(y)$，则 $f(-x)=1-A$、$f(-y)=1-B$。由完整线 $(-x,-y,w)$ 和 $(-x,y,z)$ 可得 $f(w)=A+B-1$、$f(z)=A-B$。于是完整线 $(w,z,x)$ 的标量和为 $-1$，矛盾。故最多另有一种斜率，且它与前两对中的两点共线。总行数至多 $5\cdot2=10$。各情形均与至少十二行矛盾。
\end{proof}

小行数还用到两个平面零和事实。九个斜率中必有三行之和为零。若没有同一点出现三次，就有至少五个不同点，可从中取一条完整线。五个斜率中必有非空子集之和为零。为证后者，在群代数 $\F[X,Y]/(X^3-1,Y^3-1)$ 中，每个 $1-X^aY^b$ 都属于理想 $(X-1,Y-1)$。该理想的五次方为零，因为两个生成元的三次方都为零。若五项中没有非空零和子集，则五个因子的乘积在 $X^0Y^0$ 处的系数为一，矛盾。这两个事实分别是经典平面零和结果在三元域的情形\cite{reiher,olson}。

\newpage
\section{统一上界与秩二取等构造}
令 $F(m)=m-\max\{1,\lfloor m/3\rfloor\}$，$m>0$。把域元素 $0,1,-1$ 提升为对应的整数。对大小为 $r$ 的零斜率和块，至多用 $r-1$ 次行加法就能把块内的电荷集中在一行。选 $u_j\ne0$ 的行作支点，对其余满足 $u_i\ne0$ 的行作 $R_i\gets R_i-(u_i/u_j)R_j$，即可使 $u_i$ 变为零。块的两种电荷总和都保持，最后只有支点可能有电荷，其斜率为零、标量为 $\sum_B s_i$。块内若原本 $u$ 全零，则无须操作。

先证明更一般的断言。若总斜率为零、标量总和 $t\ne0$，就能以至多 $F(m)$ 步将所有斜率电荷清零，并使最终标量的整数和 $D$ 属于 $\{\tilde t,-2\tilde t\}$，其中 $\tilde t$ 为 $t$ 的整数提升。$m\le5$ 时取整个集合为一块。$6\le m\le8$ 时从五行中取一个斜率和为零的块，其非空补集为第二块。$9\le m\le11$ 时先取零和三行，再把余下部分分为两块。分别集中各块，总代价至多为 $F(m)$。此时非零标量至多三项，故 $|D|\le3$，又有 $D\equiv t\pmod3$，断言成立。

当 $m\ge12$ 时，用引理~\ref{short} 取大小为 $r\le3$ 的块。其标量和 $\sigma$ 是 $0$ 或 $-t$，余下行的标量总和 $t-\sigma$ 仍非零。对余下行用归纳假设。若 $\sigma=0$，所得整数和已符合要求。若 $\sigma=-t$，余下行的标量总和为 $-t$，其最终整数和为 $-\tilde t$ 或 $2\tilde t$。加回本块的 $-\tilde t$，恰得 $-2\tilde t$ 或 $\tilde t$。去掉至多三行使 $\lfloor m/3\rfloor$ 至多减少一，故
\[
 r-1+F(m-r)\le F(m).
\]

对归一化输入取 $t=1$。先排零电荷，再每次按 $(-1,1)$ 排一对，最后余下 $(1)$ 或 $(-1,-1)$。所有真前缀的标量和都避开一，故同一行序对所有参数通过检验。

为证下界，令 $m=n+3$，取全体 $u_i=1$、$b=(1,0,\ldots,0)$，斜率行依次为
\[
 (1,1),\quad(-1,1),\quad(-n,1),\quad
 \underbrace{(1,0),\ldots,(1,0)}_{n\text{ 行}}.
\]
总斜率为零，前两行线性无关。初始斜率均非零，任意两行的斜率也不互为相反数。通过检验时，每行的斜率电荷必为零。若真前缀的斜率和非零，其仿射前缀和便能取到一。故每个真前缀的斜率和为零。再结合全体斜率和为零，对相邻前缀作差即得结论。随免费置换追踪行标号，把每次行加法涉及的两个标号连边。每个连通分量的斜率总和守恒，因此每个分量至少含三条初始行。若用了 $k$ 次行加法，至少有 $m-k$ 个分量，故 $k\ge m-\lfloor m/3\rfloor$。两行时取 $u=(1,1)$、$b=(1,0)$ 和一对相反的非零斜率，至少需要一步。一行时没有非空真前缀。定理~\ref{th} 得证。

\smallskip\noindent\textbf{复算与形式化}\quad
\href{https://github.com/SHIN09-git/integer-mult-bounds/tree/codex/ternary-guard-affine-quotient/research/monomial-guard-affine-quotient}{研究分支}收录精确枚举脚本和 Lean 4.24.0 证明。主定理 \texttt{worst\_case\_eq} 沿用前序工作的行操作代价 \texttt{Ternary.C}。定理 \texttt{rank\_two\_greatest} 给出秩二输入的最大值。
\begin{thebibliography}{2}\footnotesize
\bibitem{olson} J.~E. Olson, \href{https://doi.org/10.1016/0022-314X(69)90021-3}{A combinatorial problem on finite Abelian groups, I}, J. Number Theory (1969).
\bibitem{reiher} C. Reiher, \href{https://www.math.uni-hamburg.de/en/personen/reiher/artikel/Kemnitz.pdf}{On Kemnitz' conjecture concerning lattice points in the plane}, Ramanujan J. (2007).
\end{thebibliography}
\end{document}
